\chapter{可观测量与测量}
\label{ch:observables}

\section{离散谱与投影测量}

自旋分析器的每个出口既有读数，又能制备一个确定的自旋态。能量测量也给出读数，
但两个不同的原子态可能有相同的能量，仪器只亮起同一盏指示灯。
要把自旋实验的方法推广到能量、位置等可观测量，需要先确定哪些态对应同一读数，
以及筛出这个读数后，再测别的量会得到什么。谱投影把每个读数与相应的整组态联系起来。

\subsection{态空间、本征值与谱投影}
以两个出口为例，内积 $\braket{\phi}{\psi} = \phi_+\conj\psi_+ + \phi_-\conj\psi_-$ 可以由两列系数逐项相乘求和。
增加出口数目时，仍然可以如此取分量、算振幅。本书的内积对右侧线性、对左侧共轭线性。
态空间 $\HH$ 是复 Hilbert 空间，即带有内积、并且在范数 $\norm\psi = \sqrt{\braket{\psi}{\psi}}$ 下完备的复向量空间。%
\footnote{有限维内积空间自动完备。\ssecref{ss:qm-regularity}用波包的逼近说明，在无限维时为什么需要这个条件。}
右矢 $\ket\psi$ 是 $\HH$ 中的向量。给定 $\ket\phi$，对应的左矢 $\bra\phi$ 是把每个 $\ket\psi$ 映为复数 $\braket{\phi}{\psi}$ 的线性映射。

同一个自旋态可以用它在 $z$ 基中的两个分量描述，也可以用它在 $x$ 基中的两个分量描述。
选定一组正交归一基，用这组基中的坐标描述态、用相应的矩阵描述算符，就选定了一种\emph{表象}(representation，见\dfnref{dfn:qm-representation})。
在有限维空间中，选取基 $\{\ket n\}$ 后，$\psi_n = \braket{n}{\psi}$ 是态的第 $n$ 个坐标。由这些坐标可重建态矢，
\begin{equation}
    \ket{\psi} = \sum_n\ket{n}\braket{n}{\psi},
\end{equation}
同一表象中，算符 $\op A$ 的矩阵元为 $A_{mn} = \mel{m}{\op A}{n}$。用这个矩阵及两个态的坐标可以计算
\begin{equation}
    \mel{\phi}{\op{A}}{\psi} = \sum_{mn}\braket{\phi}{m}\mel{m}{\op{A}}{n}\braket{n}{\psi}.
\end{equation}
第一章的矩阵给每一列系数指定一列新的系数。一般地，线性算符 $\op{A}$ 把线性子空间 $\Dom(\op A) \subseteq \HH$ 中的向量映入 $\HH$，并保持线性组合。
这里 $\Dom(\op A)$ 称为定义域。例如对一个波函数求导，导数能否作为态空间中的向量，需要另行检查。
在有限维空间里，可以让矩阵作用于所有向量。

图\ref{fig:sg-filters}(a)中，筛出的 $+z$ 束再进入同向分析器，每次都读到 $+\hbar/2$。若改为保留 $-z$ 束，同向复测就总读到 $-\hbar/2$。
第一章用这两个读数分别给两个出口的投影加权，所得自旋算符对两种制备态的作用为
\[
    \op S_z\ket{+z} = \frac{\hbar}{2}\ket{+z},\qquad
    \op S_z\ket{-z} = -\frac{\hbar}{2}\ket{-z}.
\]
在这两式中，算符把态矢乘以该态的确定读数。这两个态矢构成一组基，由线性便可算出 $\op S_z$ 对任意叠加态的作用。
例如，对 $\ket{+x}$ 测 $S_z$ 时，两种读数等可能，平均值为零，单次读数仍是 $\pm\hbar/2$ 中的一个。

假设~\ref{psl:measurement}把这种对应推广到一般可观测量。在有限维空间中，用自伴算符 $\op A$ 表示可观测量，它的本征值给出理想测量的可能读数。
自伴算符在正交归一基中的矩阵是 Hermite 矩阵，满足 $\op A\adjoint = \op A$。有限维谱定理保证它的本征值全为实数，且可以选取一组正交归一的本征向量作基。
在本征值为 $a$ 的子空间里，$\op A$ 的作用就是乘以 $a$。因此可以先把一个态投影到各本征子空间，再将每个投影分量乘以相应的本征值并相加。\secref{sec:selfadjoint}再讨论无限维情形的自伴条件。

\begin{definition}{本征值、简并与谱投影}{qm-projector}
    先考虑有限维的自伴算符 $\op A$，即在正交归一基中由 Hermite 矩阵表示的算符。若非零向量 $\ket a$ 满足
    \begin{equation}
        \op A\ket a = a\ket a,
    \end{equation}
    称 $a$ 为本征值(eigenvalue)，归一化的本征向量称为本征态(eigenstate)。在 $a$ 对应的本征子空间中选一组正交归一基 $\ket{a,\mu}$。
    每个基矢都给出同一读数 $a$，$\mu = 1,\ldots,d_a$ 用来区分这些基矢。
    $d_a > 1$ 时称该本征值简并(degenerate)。
    本征值 $a$ 对应的谱投影(spectral projection)为
    \begin{equation}
        \op{\Pi}_a = \sum_{\mu=1}^{d_a}\ketbra{a,\mu}{a,\mu}.
    \end{equation}
    $\op{\Pi}_a\ket{\psi}$是$\ket{\psi}$中具有测量值$a$的分量。
\end{definition}

\begin{corollary}
    谱投影自伴$\op{\Pi}_a\adjoint = \op{\Pi}_a$，且投影的投影等于自身$\op{\Pi}_a^2 = \op{\Pi}_a$。
    不同本征值的投影正交$\op{\Pi}_a\op{\Pi}_b = 0$。
\end{corollary}

\begin{corollary}[有限维自伴算符的谱分解]
    有限维自伴算符 $\op A$ 的谱(spectrum)由其所有本征值组成。
    有限维 Hermite 矩阵的谱定理保证这些基矢张成全空间，因此
    \begin{subequations}
        \begin{align}
            \one &= \sum_a \op\Pi_a,\\
            \op A &= \sum_a a\op\Pi_a.
        \end{align}
    \end{subequations}
\end{corollary}

无限维自伴算符若有完备的正交归一本征基，也能按同样的谱投影展开。\secref{sec:selfadjoint}给出展开成立的条件。

\subsection{理想测量的读数与测后态}

第一章用投影筛出一路自旋，再用保留分量的范数平方算出通过率。若仪器只能分辨本征值 $a$，这一路就包含属于 $a$ 的整个本征子空间。

\begin{postulate}{可观测量与理想投影测量}{measurement}
    可观测量(observable)由自伴算符 $\op A$ 表示。离散谱情形下，测量的可能读数是 $\op A$ 的本征值。
    以 $A$ 记数值读数。对归一化纯态 $\ket\psi$，测得本征值 $a$ 的概率为
    \begin{equation}
        p_a = \P(A = a) = \mel{\psi}{\op{\Pi}_a}{\psi},
    \end{equation}
    若 $p_a > 0$，从理想 Lüders 测量中选出读数为 $a$ 的系统，它们的测后态为
    \begin{equation}
        \ket{\psi_a} = \frac{\op\Pi_a\ket\psi}{\sqrt{p_a}}.
        \label{eq:luders-update}
    \end{equation}
\end{postulate}

% 测后归一化来自
% $\|\op{\Pi}_a\ket{\psi}\|^2=\mel{\psi}{\op{\Pi}_a\adjoint \op{\Pi}_a}{\psi}=p_a$。
\begin{corollary}
    由于 $\op\Pi_a\ket{\psi_a} = \ket{\psi_a}$，再次测量测后态 $\ket{\psi_a}$ 仍得到$a$。
\end{corollary}

\subsectionstar{从测量概率的一致性到 Born 规则}
\label{ss:qm-born-consistency}
各出口的概率之和等于一。还要加上什么条件，才能推出每个出口的概率都等于相应振幅的模平方？

设归一化纯态在三个出口基矢上的分量分别为 $c_1,c_2,c_3$。记 $s_i = \abs{c_i}^2$，
则 $s_1 + s_2 + s_3 = 1$ 是向量长度的几何恒等式。先把出口概率写成未知的 $f(s_i)$。
这一步假定，在任一制备和任一分析基中，每个出口的概率只依赖自己分量的大小，而且遵循同一个函数。$\sqrt{s}$ 和 $s^2$ 都是只依赖 $s$ 的函数，还须检验它们能否使每组出口的概率之和都等于一。
若把粒子制备在第一个出口对应的基态，它就必定从第一个出口出来。三个出口的分量依次为 $1,0,0$，故 $f(1) = 1$、$f(0) = 0$。

保持第三个出口固定，只改变前两个出口的分析基。令第一个新基矢沿态在这二维子空间内的投影方向，第二个分量便为零。
转动前的三个模平方为 $(x,y,1 - x - y)$，转动后为 $(x + y,0,1 - x - y)$。对两组互斥读数，总概率都应为一。两式相减得到
\[
    f(x) + f(y) = f(x + y),\qquad x,y \geq 0,\quad x + y \leq 1 .
\]
取 $n$ 个相等的数 $1/n$ 反复用此式相加，得到 $f(1/n) = 1/n$，继而 $f(m/n) = m/n$。
又因概率非负，$f(y) - f(x) = f(y - x) \geq 0$（$y \geq x$）。用上下两侧的有理数夹住任意 $s \in [0,1]$，便有 $f(s) = s$。
所以在这些条件下，出口概率恰好是 $\abs{c_i}^2$。

为什么不能把 $\abs{c_i}^4$ 除以所有出口的四次方之和，重新归一化？例如三个模平方原为 $(1/2,1/2,0)$，
只在后两个出口之间转动分析基，使它们变为 $(1/2,1/4,1/4)$。按这个候选规则，第一个出口的概率会从 $1/2$ 升到 $2/3$，
尽管它的投影和入射态均未改变。这正说明同一出口在不同完整测量中给出相同概率，是比单次测量的总概率为一更强的要求。

能否连一开始的函数形式 $f(s_i)$ 也不预设？固定一种制备，定义从正交投影到 $[0,1]$ 的函数 $w$，给每个投影 $\op\Pi$ 指定一个概率 $w(\op\Pi)$，
并要求同一投影在任何包含它的完整测量中都取相同的概率。这项物理假设称为投影概率的非情境性。
再要求 $w(\one) = 1$，并且合并两个互相正交的出口时，概率按互斥事件相加，即
\begin{equation}\label{eq:qm-additive-projection-probability}
    w(\op\Pi + \op Q) = w(\op\Pi) + w(\op Q),\qquad \op\Pi\op Q = 0 .
\end{equation}
在维数至少为三的有限维 Hilbert 空间中，Gleason 定理\cite{gleason1957}说明，
满足这些要求的每种概率分配都对应唯一的半正定、迹为一的算符 $\op\rho$，使得
\begin{equation}\label{eq:qm-gleason-probability}
    w(\op\Pi) = \tr(\op\rho\,\op\Pi).
\end{equation}
这个算符称为密度算符。定理对实、复 Hilbert 空间均成立，因此采用复振幅仍是另一项物理选择。

现在把制备与归一化态矢 $\ket\psi$ 联系起来，要求检验该态的一维投影 $\op\Pi_\psi = \ketbra{\psi}{\psi}$ 必定通过，
即 $w(\op\Pi_\psi) = 1$。于是 $\mel{\psi}{\op\rho}{\psi} = 1$。由于 $\op\rho$ 的本征值非负且总和为一，
只有当 $\ket\psi$ 属于本征值 $1$ 的子空间时，期望值才能达到一。其余本征值都为零，故 $\op\rho = \op\Pi_\psi$。
对任一归一化分析态 $\ket a$，这就给出
\[
    w(\ketbra{a}{a}) = \tr\bigp{\ketbra{\psi}{\psi}\,\ketbra{a}{a}}
    = \abs{\braket{a}{\psi}}^2 .
\]

为什么要限定维数？在二维空间中，每个一维投影只有一个正交补。以 $\ket{+z}$ 的制备为例，
若在方向 $\bm n$ 的自旋分析器上试取正端概率 $(1 + n_z^3)/2$，其中 $n_z$ 是该方向的 $z$ 分量，
反向出口的概率就是 $(1 - n_z^3)/2$。这组连续的概率在每一对出口上都归一，并且在 $+z$ 方向给出确定读数。
但将分析器转过 $60^\circ$ 时，它预言正端通过率为 $9/16$，第一章的 Born 规则给出 $3/4$。
仅对这个二维系统的投影测量施加上述一致性条件，仍容许这两种不同的预言。

\subsection{测量的期望值与标准差}
\label{ss:measurement-moments}

在同样制备的$N$个系统上各测量一次$\op{A}$，记录读数$a$出现的次数$N_a$。
相对频率$N_a/N$逼近$p_a$。用每个读数乘以它的相对频率后求和，所得样本平均逼近期望值。
若读数的方差有限，增加独立制备的次数会减小样本平均的标准差。单次读数仍服从同一制备态的分布，仪器的分辨率由标定给出。

\begin{definition}{测量的期望值与标准差}{qm-measurement-moments}
    在离散读数的分布 $p_a$ 中，期望值(expectation value)和方差(variance)分别为
    \begin{subequations}\label{eq:qm-measurement-moments}
        \begin{align}
            \ave{\op{A}} &= \sum_a a p_a,\label{eq:qm-measurement-mean}\\
            (\Delta A)^2 &= \sum_a (a - \ave{\op{A}})^2p_a .\label{eq:qm-measurement-variance}
        \end{align}
    \end{subequations}
    标准差(standard deviation) $\Delta A$ 为方差的非负平方根。
    对连续读数，相应求和换作关于测量概率的积分；这些量要求相应的矩存在。
\end{definition}

对有限维系统，用读数对应的谱投影可以直接从态计算
\begin{subequations}
    \begin{alignat}{1}
        \ave{\op{A}} &= \sum_a ap_a = \sum_a a\mel{\psi}{\op{\Pi}_a}{\psi} = \mel{\psi}{\op{A}}{\psi},\\
        (\Delta A)^2 &= \sum_a(a - \ave{\op{A}})^2p_a = \ave{\op{A}^2} - \ave{\op{A}}^2.
    \end{alignat}
\end{subequations}
对无界的可观测量，仍可按读数的概率求均值。\ssecref{ss:qm-spectral-functions}用撤去硬墙后的能量测量说明，何时还能把这个均值写成态空间中两个向量的内积。
\begin{example}{自旋涨落与平均值的精度}{qm-spin-sampling}
    把 $N$ 个自旋 $1/2$ 粒子独立制备在 $\ket{+z}$ 态，对每个粒子测一次 $\op S_x$。求单次读数的均值和方差，以及 $N$ 次读数的平均值的标准差。

    \par\smallskip\textbf{解}\quad 两个读数 $\pm\hbar/2$ 的概率各为 $1/2$，故
    \begin{equation}
        \ave{\op{S}_x} = 0,
        \qquad
        (\Delta S_x)^2 = \frac{\hbar^2}{4}.
    \end{equation}
    两个读数大小相等、符号相反，因而平均值为零。
    记第 $r$ 次测量的读数为 $s_x^{(r)}$。对 $N$ 次独立制备，读数的样本平均 $\bar s_x$ 的方差为
    \begin{equation}
        \Var(\bar s_x)
        = \frac1{N^2}\sum_{r=1}^N\Var(s_x^{(r)})
        = \frac{\hbar^2}{4N}.
    \end{equation}
    这里独立性使不同次测量之间的协方差为零。
    样本平均的标准差为 $\hbar/(2\sqrt N)$，单次读数的标准差仍是 $\hbar/2$。
\end{example}

\subsection{仪器怎样记录测量结果}
\label{ss:measurement-record}

同一读数可以对应几个正交态。若仪器能区分这些态，它们就会留下不同的指针记录。要判断测后哪些分量还能干涉，须同时跟踪粒子和指针的状态。设指针初态为 $\ket{D_0}$，读数 $a$ 对应指针态 $\ket{D_a}$。
例如粒子有两个出口 $+,-$，指针有两种读数 $\leftLabel,\rightLabel$。把粒子出口作行、指针读数作列，联合制备的四个振幅就排成下表。
\[
   \begin{array}{c|cc}
        & \leftLabel & \rightLabel\\ \hline
      + & C_{+\leftLabel} & C_{+\rightLabel}\\
      -& C_{-\leftLabel} & C_{-\rightLabel}
   \end{array}
\]
读到 $(+,\rightLabel)$ 的概率为 $\abs{C_{+\rightLabel}}^2$。独立制备时，联合振幅是粒子振幅与指针振幅的乘积，如 $C_{+\rightLabel} = c_+d_{\rightLabel}$。两者耦合后，四个振幅未必还能这样分解。
\begin{definition}{张量积与乘积态}{qm-tensor-product}
    设粒子和仪器的态空间分别为 $\HH_{\syslabel}$ 和 $\HH_{\detlabel}$。若两边独立制备在态 $\ket{a}_{\syslabel}$ 和 $\ket{b}_{\detlabel}$，联合态称为乘积态(product state)，也称直积态。约定粒子在前、仪器在后，写作
    \begin{equation}
        \ket{a,b} = \ket{a}_{\syslabel} \otimes \ket{b}_{\detlabel} .
        \label{eq:product-basis-ket}
    \end{equation}
    $\otimes$ 表示取两个态矢的张量积(tensor product)。在两侧各选一组正交归一基，让 $a,b$ 遍历各自的基标签，所有成对的乘积态就构成联合空间 $\HH_{\syslabel} \otimes \HH_{\detlabel}$ 的一组正交归一基，称为乘积基。
    若两空间维数为 $d_{\syslabel},d_{\detlabel}$，就有 $d_{\syslabel} d_{\detlabel}$ 个乘积基矢。无限维时还要包含有限叠加在范数下的极限。张量积对两侧分别线性，例如
    \[
        (c_1\ket{s_1} + c_2\ket{s_2}) \otimes \ket{D}
        = c_1\ket{s_1} \otimes \ket{D} + c_2\ket{s_2} \otimes \ket{D}.
    \]
    两个乘积态的内积等于两侧内积之积。例如，乘积基矢的内积为 $\braket{a',b'}{a,b} = \braket{a'}{a}\braket{b'}{b}$。
\end{definition}
例如，表中 $+$ 行、$\rightLabel$ 列对应 $\ket{+,\rightLabel}$，表示粒子在 $+$ 出口而指针指向 $\rightLabel$。联合态 $\ket{\Psi}$ 在这个乘积基矢上的分量就是 $C_{+\rightLabel} = \braket{+,\rightLabel}{\Psi}$。
有些教材把式\eqref{eq:product-basis-ket}中的 $\ket{a,b}$ 连写为 $\ket{ab}$。读这种简写时，要先知道标签怎样分开、各属于哪一侧。本书把两个因子的标签合写在一个右矢内时，用逗号隔开。

如果只在仪器一侧施加算符 $\op{A}_{\detlabel}$，粒子一侧的操作就是单位算符 $\one_{\syslabel}$。由此可在联合空间上定义
\[
    (\one_{\syslabel} \otimes \op{A}_{\detlabel})(\ket{s} \otimes \ket{D}) = \ket{s} \otimes (\op{A}_{\detlabel}\ket{D}).
\]
对多个乘积态的叠加，按线性把这一式用到每一项。只写 $\op{A}_{\detlabel}$ 作用在联合态上时，默认它代表 $\one_{\syslabel} \otimes \op{A}_{\detlabel}$。
只作用于粒子的算符则写成 $\op{A}_{\syslabel} \otimes \one_{\detlabel}$。测量还要求仪器的输出依赖粒子的输入，
必须让两者耦合，才能从同一个初始指针态产生不同读数。

取一种在记录读数时保留各本征态的理想测量，其耦合可写成
\begin{equation}
    \op{U}\bigp{\ket{a,\mu} \otimes \ket{D_0}}
    = \ket{a,\mu} \otimes \ket{D_a}.
\end{equation}
这里$\mu$标记同一读数内仪器没有区分的自由度。
不同读数的指针态应正交，$\braket{D_a}{D_b} = \delta_{ab}$，
从而理想读出可以把它们完全区分。

将初态在 $\ket{a,\mu}$ 基中展开，由 $\op{U}$ 的线性和谱投影的定义，
\begin{equation}
    \op{U}\bigp{\ket{\psi} \otimes \ket{D_0}}
    = \sum_a \op{\Pi}_a\ket{\psi} \otimes \ket{D_a}.
\end{equation}
不同的输入基矢经耦合后仍彼此正交且归一化，因此耦合保持初态的范数。例如，两态都有非零振幅且产生不同的正交指针态时，联合态包含
$c_1\ket{1} \otimes \ket{D_1} + c_2\ket{2} \otimes \ket{D_2}$。
若能写成一个乘积态，沿 $\ket{1}$ 和 $\ket{2}$ 取分量得到的两个指针向量就必须成比例。这里的指针向量正交，因此联合态不可分，称为纠缠态。
在耦合后的联合态中，所有本征值为 $a$ 的分量都对应同一指针态 $\ket{D_a}$。按 Born 规则，指针显示 $a$ 的概率为 $p_a = \norm{\op\Pi_a\ket\psi}^2$。
选出指针显示 $a$ 的那些粒子，把保留下来的分量 $\op\Pi_a\ket\psi$ 除以 $\sqrt{p_a}$，便得到式\eqref{eq:luders-update}的测后态。同一读数下各分量的相对相位得以保留。

\begin{example}{简并读数的测后态}{qm-degenerate-measurement}
    在正交归一基中，$\op{A} = a(\ketbra{1}{1} + \ketbra{2}{2}) + b\ketbra{3}{3}$，其中 $a \ne b$。初态为归一化的 $c_1\ket{1} + c_2\ket{2} + c_3\ket{3}$。
    对 $\op A$ 作理想投影测量，求得到 $a$ 的概率与相应的测后态。若仪器还分辨了标签 $1,2$ 再丢弃标签，怎样用后续测量区分这两种操作？

    \par\smallskip\textbf{解}\quad 测得 $a$ 的概率是 $\abs{c_1}^2 + \abs{c_2}^2$。若此概率非零，代入式\eqref{eq:luders-update}得到
    \begin{equation}
        \ket{\psi_a} =
        \frac{c_1\ket{1} + c_2\ket{2}}{\sqrt{\abs{c_1}^2 + \abs{c_2}^2}}.
    \end{equation}
    两个分量之间的相对相位仍被保留。
    若仪器还区分了标签1与2，再把标签丢掉，
    得到的则是这两个态的统计混合。
    可以实际区分这两种操作。取 $c_1 = c_2 \ne 0$，保留读数 $a$ 后，用分析态
    $(\ket{1} + \ket{2})/\sqrt2$ 筛选。上式给出的态必定通过；
    若两个标签已被分别记录，$\ket{1}$ 束和 $\ket{2}$ 束对此筛选各有一半通过率，
    混在一起仍只有一半通过。
\end{example}

\begin{remark}[相同读数下的仪器分辨能力]
    简并本征值 $a$ 标记整个本征子空间。理想 Lüders 测量保留属于 $a$ 的各分量之间的相对相位。
    若仪器把内部标签也记录下来，只向实验者显示 $a$，随后干涉的计数仍会受到这份记录的影响。
    \secref{sec:povm}将处理仪器扰动和有限分辨率。
\end{remark}

\section{连续谱}
\label{sec:continuous}
自旋分析器的两个出口各有确定的读数；换成位置探测器，一个像素收到的粒子却可能落在一整段长度内。
若把像素做得越来越窄，用什么量还能比较同一制备下的两次位置测量？设各像素宽 $\Delta x$，第 $j$ 个收集到 $N_j$ 次，总数为 $N$。
当计数充分时，$N_j/N$ 估计这一格的概率；水平轴若按长度画，每根直方柱的高度应取
$N_j/(N\Delta x)$，这样柱的面积才等于该格的概率。若概率密度在一格内近似常量，
把这一格平分后，每个小格的概率约减半，直方柱高度近似保持原值。

\subsection{位置算符与波函数}
设一个无自旋粒子在整条直线上运动，表示其位置这一可观测量的算符称为\emph{位置算符}，记作 $\op X$。
自旋的离散本征值标记了分析器的各个出口。位置的可能读数连续变化，我们同样用实数 $x$ 标记理想的精确位置态 $\ket{x}$，写作
\begin{equation}
    \op X\ket{x} = x\ket{x}.
\end{equation}
这样的 $\ket x$ 称为广义本征态。它们可以像离散基矢一样用来展开态。离散情形用 Kronecker 符号在求和中选出一个分量：
$\braket{n}{n'} = \delta_{nn'}$。
连续情形用 Dirac $\delta$ 分布在积分中取出指定位置的函数值。设 $f$ 为光滑试验函数，则
\begin{equation}\label{eq:qm-delta-sifting}
    \sum_j\delta_{ij}c_j = c_i,
    \quad \longleftrightarrow \quad
    \int_{\RR}\delta(x - x')f(x')\d x' = f(x).
\end{equation}
因为积分的结果与 $f(x)$ 量纲相同，$\delta(x - x')$ 的量纲是长度的倒数。位置基的归一化关系为
\begin{equation}
    \label{eq:qm-position-basis}
    \braket{x}{x'} = \delta(x - x').
\end{equation}
对应的，
与离散基的完备关系
$\textstyle \sum_n \ketbra{n}{n} = \one$
类似，位置基的完备关系为
\begin{equation}
    \label{eq:qm-position-completeness}
    \int_\RR\d x\,\ketbra{x}{x} = \one.
\end{equation}

\begin{definition}{位置表象的波函数}{qm-wave-function}
    态 $\ket{\psi}$ 在位置基中的分量称为位置波函数(wave function)，记作
    \begin{equation}\label{eq:qm-position-wavefunction}
        \psi(x) = \braket{x}{\psi}.
    \end{equation}
\end{definition}

\begin{corollary}[波函数的归一化]
    若 $\ket\psi$ 已归一化，在内积中插入式\eqref{eq:qm-position-completeness}，得到
    \begin{equation}\label{eq:qm-position-normalization}
        1 = \braket{\psi}{\psi}
        = \int_\RR \braket{\psi}{x}\braket{x}{\psi}\d x
        = \int_\RR \abs{\psi(x)}^2\d x .
    \end{equation}
\end{corollary}

在位置表象中，态用平方可积函数表示，所用的 Hilbert 空间为 $L^2(\RR,\d x)$。
用 $X$ 表示探测器给出的位置读数。按 Born 规则，覆盖区间 $[a,b]$ 的像素探测到粒子的概率为
\begin{equation}\label{eq:qm-position-interval-probability}
    \P(a \leq X \leq b) = \int_a^b\abs{\psi(x)}^2\d x .
\end{equation}
因此 $\abs{\psi(x)}^2$ 是位置概率密度，单位是长度的倒数。例如，若位置均匀分布在一段 $\qty{1}{cm}$ 的区间内，落在其中一个 $\qty{1}{mm}$ 像素内的概率为 $0.1$。
位置读数 $X$ 落在区间 $[a,b]$ 内，对应的投影为
\begin{equation}\label{eq:qm-position-spectral-projection}
    \op{\Pi}_X([a,b]) = \int_a^b\ketbra{x}{x}\d x .
\end{equation}
这里 $\op{\Pi}_X(\RR) = \one$。互不相交的读数区间对应正交投影，它们相加就给出区间并集的投影。
% 下标 $X$ 指明测量的是位置。换成只接受某段动量的探测器，所求的就是相应动量区间的谱投影在制备态中的期望值。
% 这个期望值可以在任意表象中计算。选用动量基，就能直接对接受区间内的动量概率密度积分。

位置算符 $\op X$ 作用于一般态 $\ket\psi$ 后得到的右矢 $\op X\ket\psi$ 是什么？
类似得到波函数 $\braket{x}{\psi} = \psi(x)$ 的处理，我们考虑 $\op X\ket\psi$ 在位置基中的分量 $\mel{x}{\op X}{\psi}$。
由本征方程 $\op X\ket x = x\ket x$ 和 $\op X$ 的自伴性可得
\begin{equation}
    \mel{x}{\op X}{\psi} = x\braket{x}{\psi} = x\psi(x).
\end{equation}
每个位置分量都乘以相应的本征值。
这要求 $\ket\psi$ 属于 $\op X$ 的定义域，即 $x\psi(x)$ 仍平方可积。两个波函数若只在零测集上取值不同，就代表同一向量。这里的函数等式也按几乎处处相等来理解。


\subsection{动量算符}
\label{ss:momentum-position}
表示粒子沿 $x$ 轴的动量这一可观测量的算符称为\emph{动量算符}，记作 $\op P$。
用 $p$ 标记可能的动量读数，相应的广义本征态 $\ket{p}$ 满足
\begin{equation}
    \op P\ket{p} = p\ket{p}.
\end{equation}
要求 $\op P\ket\psi$ 的位置分量，先要知道动量确定的态 $\ket p$ 在位置基中有哪些分量 $\braket{x}{p}$。

在一维自由空间，可以把波包分解成平面波。以下推导先用 Gauss 波包一类光滑、迅速衰减的波包，\ssecref{ss:qm-regularity}说明如何放宽这个条件。
用带符号的波数 $k$ 标记每个分量，$k$ 的正负区分传播方向，非零波数对应的波长为 $\lambda = 2\pi/\abs{k}$。
采用对称的 Fourier 变换约定，波数态的位置波函数为 $\braket{x}{k} = \e^{\i kx}/\sqrt{2\pi}$，这是为了保证归一化
\begin{equation}
    \braket{k}{k'} = \int_\RR\d x\,\braket{k}{x}\braket{x}{k'}
    = \frac{1}{2\pi}\int_\RR\d x\,\e^{-\i (k-k')x}
    = \delta(k - k').
\end{equation}
要从 $\psi(x)$ 中取出这个分量，就与它的复共轭 $\braket{k}{x}$ 做积分，得到波数表象的振幅
\begin{equation}
    \braket{k}{\psi} = \int_\RR\d x\,\braket{k}{x}\braket{x}{\psi} = \frac{1}{\sqrt{2\pi}}\int_\RR\d x\,\psi(x)\e^{-\i kx}.
\end{equation}
怎样知道波数为 $k$ 的分量对应动量 $p$？在\secref{sec:interference}的电子衍射实验中，加速电压给出电子的动量大小，衍射峰的位置与已知晶面间距给出波长。
本节以这两种测量所检验的 de Broglie 关系 $p = \hbar k$ 为物理出发点，
为了满足归一化条件 $\braket{p}{p'} = \delta(p - p')$，动量本征态的位置波函数确定为
\begin{equation}
    \braket{x}{p} = \frac{1}{\sqrt{\hbar}}\bigeval{\braket{x}{k}}_{k=p/\hbar}
    = \frac{\e^{\i px/\hbar}}{\sqrt{2\pi\hbar}}.
\end{equation}
% 三个方向的平面波归一化因子相乘，得到三维的归一化因子 $(2\pi\hbar)^{-3/2}$。
将 $\psi(x)$ 投影到这个平面波上，得到动量波函数
\begin{equation}
    \widetilde\psi(p) = \braket{p}{\psi}
    = \int_\RR\d x\,\braket{p}{x}\braket{x}{\psi}
    = \frac{1}{\sqrt{2\pi\hbar}}\int_\RR\d x\,\psi(x)\e^{-\i px/\hbar}.
\end{equation}

\paragraph{动量测量的概率}
一维位置和动量振幅的量纲分别为长度$^{-1/2}$和动量$^{-1/2}$。
把一次动量测量的随机读数记为 $P$，它落在区间 $B$ 内的概率记为 $\P(P \in B)$。投影 $\op\Pi_P(B)$ 保留这个区间内的动量分量。对归一化态，所保留分量的范数平方就是该概率，
\begin{subequations}\label{eq:qm-momentum-measurement}
    \begin{align}
        \op\Pi_P(B) &= \int_B\ketbra{p}{p}\d p,\\
        \P(P \in B) &= \mel{\psi}{\op\Pi_P(B)}{\psi} = \int_B\abs{\widetilde\psi(p)}^2\d p .
    \end{align}
\end{subequations}

\begin{example}{动量探测器的接受概率}{qm-momentum-window}
    一个自由粒子在 $\abs{p} < p_0$ 内的动量振幅为常数 $C$，其余处为零，$p_0 > 0$。
    一台探测器接受 $-p_0/2 \leq p \leq p_0/2$ 的粒子。重复制备后，通过的比例是多少？

    \par\smallskip\textbf{解}\quad 归一化要求 $2p_0\abs{C}^2 = 1$，故密度为 $1/(2p_0)$。
    接受窗口的宽度为 $p_0$，由式\eqref{eq:qm-momentum-measurement}对该窗口积分，得到 $p_0/(2p_0) = 1/2$。
    在窗口内继续缩小宽度，概率随宽度成比例减小，密度保持不变。
\end{example}

\begin{example}{位置与动量的 Gauss 波包}{qm-gaussian-wavepacket}
    粒子的位置波函数为 $\psi(x) \propto \exp[-x^2/(4\sigma^2)]$，其中 $\sigma > 0$。求位置和动量读数的概率密度，以及两种读数的标准差之积。

    \par\smallskip\textbf{解}\quad 将位置波函数代入 Fourier 积分，动量振幅正比于 $\exp(-\sigma^2p^2/\hbar^2)$。将两个振幅分别取模平方并归一化，得到
    \begin{subequations}\label{eq:qm-gaussian-conjugate-densities}
        \begin{align}
            \abs{\psi(x)}^2 &= \frac{\e^{-x^2/(2\sigma^2)}}{\sqrt{2\pi}\,\sigma},\\
            \abs{\widetilde\psi(p)}^2 &= \sqrt{\frac{2}{\pi}}\,\frac{\sigma}{\hbar} \e^{-2\sigma^2p^2/\hbar^2}.
        \end{align}
    \end{subequations}
    与标准 Gauss 密度比较指数的二次项，可读出 $\Delta X = \sigma$、$\Delta P = \hbar/(2\sigma)$。把位置分布压窄会使动量分布展宽，这一族波包的宽度乘积始终为 $\Delta X\Delta P = \hbar/2$。
\end{example}

\paragraph{动量算符在两种表象中的作用}
动量算符作用后，右矢 $\op P\ket\psi$ 的动量分量等于原振幅乘以相应的动量读数，
\begin{equation}
    \mel{p}{\op P}{\psi} = p\braket{p}{\psi} = p\widetilde\psi(p).
\end{equation}
在位置基中，把 $\ket\psi$ 展开成平面波。由于 $\p_x\braket{x}{p} = (\i p/\hbar)\braket{x}{p}$，每个平面波乘以 $p$ 等价于对它求位置导数再乘以 $-\i\hbar$。因此
\begin{equation}
    \label{eq:momentum-position-derivation}
    \begin{aligned}
        \mel{x}{\op P}{\psi}
        &= \int_\RR\d p\,\mel{x}{\op P}{p}\braket{p}{\psi}
         = \int_\RR\d p\,p\braket{x}{p}\braket{p}{\psi}\\
        &= -\i\hbar\int_\RR\d p\,\pv{\braket{x}{p}}{x}\braket{p}{\psi}
         = -\i\hbar\dd x\int_\RR\d p\,\braket{x}{p}\braket{p}{\psi} = -\i\hbar\dv{\braket{x}{\psi}}{x}.
    \end{aligned}
\end{equation}
于是动量算符在位置表象中的作用为
\begin{equation}
    \mel{x}{\op P}{\psi} = -\i\hbar\psi'(x).
    \label{eq:momentum-position}
\end{equation}

用 $\op X$ 和 $\op P$ 在位置表象中的作用规则，可以体现 $\op X\op P \neq \op P\op X$。对前面选取的光滑波包，
\begin{equation}\label{eq:canonical-position-momentum}
    \mel{x}{[\op X,\op P]}{\psi}
    = -\i\hbar x\psi'(x) + \i\hbar\bigp{x\psi(x)}'
    = \i\hbar\psi(x).
\end{equation}
对每个这样的波包，等式在几乎所有 $x$ 上成立。位置波函数确定了整个右矢，所以可简记为
\begin{equation}
    [\op X,\op P] = \i\hbar\one.
\end{equation}
这个算符关系不依赖具体表象。\chapref{ch:angular}从 Weyl 关系的指数形式作算符推导。

\section{可观测量的相容性与不确定关系}
能否换一种波包形状，使位置与动量的标准差之积小于 $\hbar/2$？先回到可以逐个计数的自旋测量。沿 $z$ 筛出的原子，重复测 $\op{S}_z$ 有确定的结果，改测 $\op{S}_x$ 却有涨落。
把同一种制备分成两批，各测一个量，两批读数的标准差能否同时任意小？
第一章的自旋投影说明哪些出口会被保留。一般地，谱投影可以用来判断哪些量能一起精确测量，态的内积则给出两个标准差应当满足的不等式。

\subsection{对易性与共同本征基}
设在同一系统上测量 $A$ 和 $B$。用读数给共同本征态加上标签 $\ket{a,b}$，表示
\begin{subequations}\label{eq:qm-joint-eigenstate}
    \begin{align}
        \op A\ket{a,b} &= a\ket{a,b},\\
        \op B\ket{a,b} &= b\ket{a,b}.
    \end{align}
\end{subequations}
这里的逗号分隔同一态的两个读数，这两个量也可以都属于同一个粒子。若这对读数仍然简并，可再加一个标签 $\mu$，写作 $\ket{a,b,\mu}$。
若能选出这样的正交归一基，就能把任意制备展开在共同基里，直接求出两种读数的联合概率。怎样事先判断有无这样的基？
\begin{theorem}{对易与共同本征基}{qm-compatible}
    有限维空间上，两个自伴算符 $\op A,\op B$ 满足 $[\op A,\op B] = 0$，当且仅当存在由它们的共同本征向量组成的正交归一基。
\end{theorem}
\begin{proof}
设两算符对易，先把 $\op A$ 对角化。在本征值 $a$ 的子空间中任取 $\ket v$，就有
\[
   \op A(\op B\ket v) = \op B(\op A\ket v) = a\,\op B\ket v .
\]
所以 $\op B$ 把这个子空间映回自身。若其中只有一个独立方向，$\op B\ket v$ 就是 $\ket v$ 的数倍；
若 $a$ 简并，$\op B$ 限制在这个子空间上仍然自伴，可以在其中再选一组 $\op B$ 的正交归一本征向量。
对每个 $a$ 都这样做，把得到的基合起来就是全空间的共同本征基。
反过来，在每个共同本征向量上，两种作用次序分别乘以 $ab$ 和 $ba$，结果相同。
若这些向量张成全空间，由线性，两种次序就在每个态上相同。
\end{proof}
\begin{remark}
    这里需要一组\emph{完备}的共同本征向量。单独找到一个共同本征态，只能说明在那一态上两种读数都确定。
    同时讨论两个可观测量时，记 $\op\Pi_a^A = \op\Pi_A(\{a\})$、$\op\Pi_b^B = \op\Pi_B(\{b\})$，分别筛出 $A = a$、$B = b$ 的分量。
    它们在共同本征基上都是只含 $0,1$ 的对角矩阵，
    故 $\op\Pi_a^A\op\Pi_b^B = \op\Pi_b^B\op\Pi_a^A$ 正好保留两个条件都满足的分量。
    对于理想投影测量，两种筛选顺序给出相同的联合概率，
    \begin{equation}
        \P(A = a,B = b) = \norm{\op\Pi_b^B\op\Pi_a^A\ket\psi}^2
        = \norm{\op\Pi_a^A\op\Pi_b^B\ket\psi}^2 .
    \end{equation}
    在无限维空间可以直接用区间筛选来表述相容性。若任取两个读数区间，相应的谱投影都对易，就称这两个可观测量\emph{强对易}。这给出了区间读数的联合谱测量，即使精确读数没有可归一化的本征向量也适用。对于无界算符，只在某个稠密定义域上算得对易子为零，还不足以推出强对易。
\end{remark}

\subsection{测量涨落与不确定关系}
现在分两批测量，在同一归一化态 $\ket\psi$ 上分别求 $A,B$ 的读数分布。按\dfnref{dfn:qm-measurement-moments}，
用大写 $\Delta A$ 表示 $A$ 的标准差，用小写 $\delta\op A$ 表示从算符中减去均值后的算符。对离散读数，
\begin{subequations}\label{eq:qm-operator-fluctuation}
    \begin{align}
        \delta\op A &= \op A - \ave{\op A}\one,\\
        \Delta A &= \sqrt{\sum_a (a - \ave{\op A})^2p_a}
                = \norm{\delta\op A\ket\psi}.
    \end{align}
\end{subequations}
最后一个等号来自把 $\ket\psi$ 按 $\op A$ 的谱投影分解后对各正交分量的范数平方求和。连续读数用积分代替求和，同理定义 $\Delta B$。
于是 $\Delta A = 0$ 恰好意味着 $\op A\ket\psi = \ave{\op A}\ket\psi$，即这一制备每次都给出相同的 $A$ 读数。

\begin{theorem}{不确定关系}{qm-uncertainty}
    设 $\op{A},\op{B}$ 自伴且归一化态 $\ket{\psi}$ 属于两者的定义域，取
    $\ket{u} = \delta \op{A}\ket{\psi}$、$\ket{v} = \delta \op{B}\ket{\psi}$。则
    $(\Delta A)^2(\Delta B)^2 \geq \abs{\braket{u}{v}}^2$，等号成立当且仅当两个涨落态线性相关。
    若算符乘积 $\op{A}\op{B}$ 和 $\op{B}\op{A}$ 在此态上也有定义，下界可写成
    \begin{equation}
        (\Delta A)^2(\Delta B)^2 \geq \abs{\braket{u}{v}}^2
            = \frac14\bigabs{\ave{[\op{A},\op{B}]}}^2
            + \frac14\bigabs{\ave{\{\delta \op{A},\delta \op{B}\}}}^2.
        \label{eq:uncertainty}
    \end{equation}
\end{theorem}
\begin{proof}
这个证明先把 Cauchy--Schwarz 不等式用于两个涨落态。由它们的定义，
有$\braket{u}{u} = (\Delta A)^2$、$\braket{v}{v} = (\Delta B)^2$。
当$\ket{v} \ne 0$时，将$\ket{u}$分解为沿$\ket{v}$和垂直于$\ket{v}$的两部分：
\begin{equation}
    \ket{u} = \ket{v}\frac{\braket{v}{u}}{\braket{v}{v}} + \ket{u_\perp},
    \qquad \braket{v}{u_\perp} = 0.
\end{equation}
两部分正交，范数平方相加，得到
\[
    \braket{u}{u}\,\braket{v}{v}
      = \abs{\braket{u}{v}}^2 + \braket{u_\perp}{u_\perp}\braket{v}{v}
      \geq \abs{\braket{u}{v}}^2 .
\]
当 $\ket v = 0$ 时也成立。
接着利用
\begin{equation}
    \delta \op{A}\delta \op{B}
    = \frac12\{\delta \op{A},\delta \op{B}\} + \frac12[\op{A},\op{B}].
\end{equation}
若两个乘积有定义，就有 $\braket{u}{v} = \mel{\psi}{\delta \op{A}\delta \op{B}}{\psi}$。
上式中反对易项的期望为实数，对易项的期望为纯虚数；模平方取两者平方之和，便得式\eqref{eq:uncertainty}。
当$\ket{v} \ne 0$时，范数不等式取等号要求$\ket{u_\perp} = 0$。
当$\ket{v} = 0$时，两涨落态已线性相关，不等式两侧均为零。
\end{proof}

\begin{corollary}[Robertson 不确定关系]
    从式\eqref{eq:uncertainty}删去非负的协方差(covariance)平方项，便有
    \begin{equation}\label{eq:qm-robertson-inequality}
        \Delta A\,\Delta B \geq \frac12\bigabs{\ave{[\op{A},\op{B}]}}.
    \end{equation}
    此式取等号要求两个涨落态线性相关，且 $\ave{\{\delta \op{A},\delta \op{B}\}} = 0$。
\end{corollary}

\begin{corollary}
    对整条实线上位置和动量的标准差都有限的波包，有
    \begin{equation}
        \Delta X \Delta P \geq \frac\hbar2.
    \end{equation}
\end{corollary}
对光滑且迅速衰减的波包，把正则对易关系\eqref{eq:canonical-position-momentum}代入 Robertson 不等式即可。对一般的有限方差波包，$x\psi$ 和弱导数 $\psi'$ 都平方可积（见\ssecref{ss:qm-regularity}）。
令 $g(x) = x\abs{\psi(x)}^2$。$g$ 和 $g' = \abs{\psi}^2 + 2x\Re(\psi\conj\psi')$ 都绝对可积，故 $g$ 在两端趋于零。对 $g'$ 在全轴上积分，利用波函数的归一化，对位置和动量的两个涨落态便有
\[
    \Im\braket{u}{v} = -\hbar\,\Re\int_\RR x\psi\conj(x)\psi'(x)\d x = \frac{\hbar}{2} .
\]
由 Cauchy--Schwarz 不等式，$\Delta X\Delta P \geq \abs{\braket{u}{v}} \geq \hbar/2$。
\begin{corollary}
    对自旋，
    \begin{equation}
        \Delta S_x \Delta S_y \geq \frac{\hbar}{2}\abs{\ave{\op{S}_z}}.
    \end{equation}
    这个下界取决于制备态，在 $\ave{\op{S}_z} = 0$ 时为零。
\end{corollary}

时间在此是外部参数，并无普适的自伴时间算符$\op{T}$满足$[\op{T},\op{H}] = \i\hbar\one$。
对不显含时间的 $\op A$，式\eqref{eq:heisenberg}给出 $\abs{\d\ave{\op A}/\nd t} = \abs{\ave{[\op H,\op A]}}/\hbar$。将它代入 $\op A$ 与 $\op H$ 的 Robertson 不等式，若均值的变化率非零，得到
\begin{equation}
    \tau_A\Delta H \geq \frac\hbar2,\qquad
    \tau_A := \frac{\Delta A}{\abs{\d\ave{\op{A}}/\nd t}},
\end{equation}
$\tau_A$ 表示按当前变化率估算，均值改变一个标准差所需的时间。

\subsection{测量仪器的误差与扰动}
\label{ss:qm-error-disturbance}
同样制备两束 $\ket{+z}$ 中子，分别用理想分析器测 $\op\sigma_x$ 和 $\op\sigma_y$，两束都给出各半的正、负读数。
即使分析器完全准确，两个标准差也都是 1。这正是式\eqref{eq:qm-robertson-inequality}讨论的制备涨落。
现在让中子依次通过两台仪器。第一台要读出 $A$，它有多准确？中子经过它以后，原先的 $B$ 受到多大扰动？
这两个问题还需要知道第一台仪器怎样与中子耦合。

沿用\ssecref{ss:measurement-record}的指针模型，记系统与指针的初态为 $\ket{\Psi} = \ket\psi\otimes\ket{D_0}$，
第一台仪器的耦合为 $\op U$，指针算符 $\op M$ 的刻度与 $A$ 相同。
为在同一个初态中比较输入与输出，记 $\op A_{\incoming} = \op A\otimes\one$、$\op B_{\incoming} = \op B\otimes\one$，并将输出算符拉回输入端，
\[
    \op M_{\outgoing} = \op U\adjoint(\one\otimes\op M)\op U,\qquad
    \op B_{\outgoing} = \op U\adjoint(\op B\otimes\one)\op U .
\]
Ozawa 用差算符在联合初态中的均方根定义测量误差 $\epsilon(A)$ 和扰动 $\eta(B)$\cite{ozawa-noise}，
\begin{subequations}
    \begin{align}
        \epsilon(A) &= \norm{(\op M_{\outgoing} - \op A_{\incoming})\ket{\Psi}},\\
        \eta(B) &= \norm{(\op B_{\outgoing} - \op B_{\incoming})\ket{\Psi}}.
    \end{align}
\end{subequations}
两项分别比较指针与目标量、耦合前后的同一个量。差算符的平方含有交叉项，实验上可借助若干种输入制备重建这些项。
在有关方差和对易子均有定义时，这组定义给出
\begin{equation}\label{eq:qm-ozawa}
    \epsilon(A)\eta(B) + \epsilon(A)\Delta B + \Delta A\,\eta(B)
        \ge \frac12\bigabs{\mel{\psi}{[\op A,\op B]}{\psi}} .
\end{equation}
其由来可从输出端看出。指针与系统的 $B$ 作用在不同因子上，故 $[\op M_{\outgoing},\op B_{\outgoing}] = 0$。
将两者各写成输入算符加上相应差算符并展开，输入对易子便等于其余三个对易子之和的负值。
对各项用三角不等式和 Robertson 不等式，再用差算符的均方根估计其标准差，即得式\eqref{eq:qm-ozawa}。

2012 年，\phantomsection\label{exp:qm-neutron-disturbance}Erhart 等用两台中子自旋分析器分别估计第一台的测量误差及其对后续测量的扰动\cite{erhart-neutron}。目标为 $\op A = \op\sigma_x$、$\op B = \op\sigma_y$。
第一台的轴可在 $xy$ 平面转动，实际沿 $\op\sigma_\phi = \op\sigma_x\cos\phi + \op\sigma_y\sin\phi$ 作理想投影测量，
第二台始终沿 $y$ 读出。记第一台的两个投影为 $\op\Pi_\pm = (\one \pm \op\sigma_\phi)/2$。
对于保留各出口本征态的仪器，输出 $s$ 对应的未归一化态是 $\op\Pi_s\ket\psi$。
在各指针出口分别取分量后，误差平方为 $\norm{(\op\sigma_\phi - \op\sigma_x)\ket\psi}^2$，
扰动平方为 $\textstyle\sum_{s=\pm}\norm{[\op\Pi_s,\op\sigma_y]\ket\psi}^2$。
用 Pauli 乘法规则，对 $\ket\psi = \ket{+z}$、$0 \le \phi \le \pi/2$ 得
\[
    \epsilon(A) = 2\sin(\phi/2),\qquad \eta(B) = \sqrt{2}\cos\phi .
\]
在 $\phi = 0$ 处，第一台准确测 $A$，误差与扰动的乘积为零，而输入对易子给出的右端为 1。
此时式\eqref{eq:qm-ozawa}左端仍有 $\Delta A\,\eta(B) = \sqrt{2}$。
实验对每个选定的 $\phi$ 分别制备 $+z,-z,+x,+y$ 四束，记录每种制备下两台分析器的四种正负组合计数。
这些计数重建出 $\ket{+z}$ 输入的误差和扰动，显示只留乘积项的下界可以被突破，三项之和与上述关系相符。

\sectionstar{算符的定义域、边界条件与谱}
\label{sec:selfadjoint}
在直线上，式\eqref{eq:momentum-position}用位置波函数的一阶导数表示动量的作用。要让结果仍是态空间中的向量，哪些波包可以求导？若粒子只能沿环运动，还要给出怎样的边界条件？这些条件决定算符的作用和可能的测量读数。

\subsectionstar{波包与算符的定义域}
\label{ss:qm-regularity}
算符的定义域 $\Dom(\op A)$ 是它允许作用的向量集合，作用结果仍在态空间 $\HH$ 中。若 $\HH$ 中的每个向量都能按范数用该域中的向量逼近，就称定义域稠密(dense)。

在整条实线上，Gauss 波包及其各阶导数在无穷远处都比任意反幂衰减得更快。具有这一性质的光滑函数称为 Schwartz 函数。
它们的 Fourier 变换仍是 Schwartz 函数。因此在式\eqref{eq:momentum-position-derivation}中，$\widetilde\psi(p)$ 和 $p\widetilde\psi(p)$ 都绝对可积，后者的模乘上一个常数就给出了求导后被积函数的可积上界，可以交换求导和积分。
在 Schwartz 波包上乘以 $x$、求导或平移，所得函数仍属于 Schwartz 空间，所以式\eqref{eq:canonical-position-momentum}中的两个算符乘积都有定义。
任何平方可积态都能按范数用 Schwartz 波包逼近。Fourier 变换保持范数，因而可以用这些波包的变换取范数极限，定义一般态的动量振幅。

\paragraph{动量的定义域与弱导数}
在动量表象中，$\op P$ 作用后得到分量 $p\widetilde\psi(p)$，所以整条实线上的动量定义域恰为
\begin{equation}\label{eq:qm-momentum-domain}
    \Dom(\op P) = \biggset{\ket\psi \in \HH}{\int_{\RR}p^2\abs{\widetilde\psi(p)}^2\d p < \infty}.
\end{equation}
在位置表象中，可以用分部积分来定义求导。若平方可积函数 $\chi$ 对每个光滑且有紧支撑的测试函数 $g$ 都满足
\begin{equation}\label{eq:qm-weak-spatial-derivative}
    \int_{\RR} g(x)\chi(x)\d x = -\int_{\RR}g'(x)\psi(x)\d x ,
\end{equation}
就称 $\chi$ 为 $\psi$ 的平方可积弱导数，记作 $\psi'$。Fourier 变换把这个导数对应到 $\i p\widetilde\psi(p)/\hbar$，所以具有平方可积弱导数的波包恰好构成式\eqref{eq:qm-momentum-domain}的动量定义域。
这也是差商 $[\psi(x-a)-\psi(x)]/a$ 在 $a \to 0$ 时按平方可积范数收敛到 $-\psi'$ 的条件，第三章将用这一变化率解释动量怎样生成平移。

\paragraph{矩形波包与动量截断}
哪些平方可积的态没有平方可积的弱导数？先看整条实线上的矩形波包。它在宽为 $L$ 的区间内振幅处处相同，直接作 Fourier 变换得
\begin{subequations}\label{eq:qm-rectangular-packet-fourier}
    \begin{align}
        \psi(x) &=
        \begin{cases}
            L^{-1/2},&\abs{x} < L/2,\\
            0,&\abs{x} \geq L/2,
        \end{cases}\\
        \widetilde\psi(p) &= \sqrt{\frac{L}{2\pi\hbar}}\frac{\sin u}{u},
            \qquad u = \frac{pL}{2\hbar}.
    \end{align}
\end{subequations}
在 $u = 0$ 处取比值的极限 $1$。位置密度的积分为 $1$，由 Fourier 变换保持范数，动量密度也已归一化。
若暂时只保留 $\abs{p} \leq \Lambda$ 的分量，记所得波包为 $\psi_\Lambda$。当 $\Lambda \to \infty$ 时，
\[
    \norm{\psi_\Lambda - \psi}^2
        = \int_{\abs{p}>\Lambda}\abs{\widetilde\psi(p)}^2\d p
        \longrightarrow 0 .
\]
若取所有具有有限动量截断的波包，截断上限可以随波包而变，所得集合在加法和数乘下封闭。让截断上限逐次趋于无穷，就得到一列在范数下越来越接近的波包。这个例子说明了为什么要考虑下述完备性。
\begin{definition}{Cauchy 列、完备性与 Hilbert 空间}{qm-hilbert}
    在赋范向量空间 $V$ 中，两个向量 $u,v$ 的距离是 $\norm{u - v}$。
    对 $V$ 中的一列向量 $(v_n)$，若对每个 $\epsilon > 0$，都存在 $N$，使得 $m,n \ge N$ 时 $\norm{v_m - v_n} < \epsilon$，就称 $(v_n)$ 为 \emph{Cauchy 列}。

    若 $V$ 中的每个 Cauchy 列 $(v_n)$ 都有极限 $v \in V$，即 $\norm{v_n - v} \to 0$，就称 $V$ 在这个范数下\emph{完备}。完备的赋范向量空间称为 \emph{Banach 空间}。

    在复内积空间中，内积诱导范数 $\norm{\psi} = \sqrt{\braket{\psi}{\psi}}$。若空间在这个范数下完备，就称它为复 \emph{Hilbert 空间}。
\end{definition}
上面的截断波包构成 Cauchy 列，极限是矩形波包。它的动量振幅在任何有限的截断之外都有尾部，因此有限动量截断的空间缺少这个极限。平方可积波函数构成的 Hilbert 空间 $L^2(\RR,\d x)$ 则包含它。
再用式\eqref{eq:qm-momentum-domain}检验动量能否作用于这个极限态。
在动量二阶矩的积分中，权重 $p^2$ 恰好抵消平方振幅分母里的 $p^2$。取 $\Lambda = 2\pi n\hbar/L$，其中 $n$ 为正整数，此时积分覆盖整数个振荡周期，
\[
    \int_{-\Lambda}^{\Lambda}p^2\abs{\widetilde\psi(p)}^2\d p
       = \frac{2\hbar}{\pi L}\int_{-\Lambda}^{\Lambda}
           \sin^2\Bigp{\frac{pL}{2\hbar}}\d p
       = \frac{2\hbar\Lambda}{\pi L}.
\]
截断态的范数趋于 $1$，重新归一化后，二阶矩仍随截断上限线性增长。自由粒子的动能为 $p^2/(2m)$，故这个矩形波包的平均动能发散。
在位置表象中也能看出原因。若把高度为 $A = L^{-1/2}$ 的一侧边缘改成宽度为 $\epsilon$ 的直线斜坡，坡上的斜率大小为 $A/\epsilon$，这一段对导数模平方积分的贡献便是 $A^2/\epsilon$。
让 $\epsilon$ 趋于零，积分趋于无穷。理想矩形的函数值在两端跳变，按分部积分求一次导数就在两端各产生一个 Dirac $\delta$ 项，因此它没有平方可积的一阶弱导数。

\paragraph{撤去硬墙后的波包}
同样只在一个有限区间内非零，波包也可以有有限的平均动能。在 $x = \pm L/2$ 放置硬墙，取阱内势能为零。把\ssecref{ss:position-well}的坐标原点移到阱中心，基态在阱内为 $\phi(x) = \sqrt{2/L}\cos(\pi x/L)$。
它在两端连续降到零，向墙外延拓为零后，跳变的是斜率。此时一阶弱导数在阱内等于余弦波的导数，在阱外为零。
它的模平方可积。在动量表象中，一阶求导对应乘以 $\i p/\hbar$。由 Fourier 变换保持范数的性质，撤墙瞬间的平均自由动能为
\[
    \int_{\RR}\frac{p^2}{2m}\abs{\widetilde\phi(p)}^2\d p
    = \frac{\hbar^2}{2m}\int_{-L/2}^{L/2}\abs{\phi'(x)}^2\d x
    = \frac{\hbar^2\pi^2}{2mL^2}.
\]
从大动量处也能看出差别。对矩形波包的 Fourier 积分作一次分部积分，端点项就给出按 $1/\abs{p}$ 衰减的振幅包络。
对阱内的余弦波作同样的分部积分，第一次的端点项为零，第二次才出现端点斜率。因此 $\widetilde\phi(p) = O(\abs{p}^{-2})$，动能积分的尾部收敛。

撤墙瞬间的平均动能等于有墙时的基态能量。有墙时本征方程只在阱内求二阶导数，阱内有 $-\phi'' = (\pi/L)^2\phi$，两端满足 $\phi(\pm L/2) = 0$。
撤墙后，自由 Hamilton 算符为 $\op H_0 = \op P^2/(2m)$。在整条实线上对零延拓求二阶导数，两端的斜率跳变各产生一个 $\delta$ 项，微分表达式 $-\hbar^2\phi''/(2m)$ 给出的是一个无法用平方可积函数表示的分布。

所以撤墙瞬间的平均动能有限，该态却超出了自由 Hamilton 算符的定义域。\ssecref{ss:qm-spectral-functions}将用能量分布的一阶矩和二阶矩区分这两个条件。

\Needspace{7\baselineskip}
上面用位置波函数的导数和动量波函数的积分描述了同一个定义域。一般地，换表象时向量、算符和定义域按同一次变换一起改写。
\begin{definition}{态空间的表象}{qm-representation}
    为态空间 $\HH$ 选定一种表象，就是选定一个容纳分量的具体 Hilbert 空间 $\mathcal K$ 和一个幺正（酉）同构 $U\colon \HH \to \mathcal K$。这意味着 $U$ 是保持内积的线性双射。
    对于态 $\ket\psi$，$U\ket\psi$ 给出它的全部分量。算符 $\op A$ 在这个空间中的表示为 $A_U = U\op A U^{-1}$，其定义域为 $\Dom(A_U) = U\Dom(\op A)$。
\end{definition}
例如，对直线上无内部自由度的粒子，位置表象取 $\mathcal K = L^2(\RR,\d x)$，$U_X\ket\psi$ 是位置波函数 $\psi(x) = \braket{x}{\psi}$。只在零测集上不同的函数代表同一向量。
动量表象取 $\mathcal K = L^2(\RR,\d p)$，$U_P\ket\psi$ 是动量波函数 $\widetilde\psi(p) = \braket{p}{\psi}$。前面从位置振幅求动量振幅的 Fourier 变换就是 $U_P U_X^{-1}$。

\subsectionstar{边界条件与自伴性}
\label{ss:qm-boundary-selfadjoint}
把长为 $L$ 的环在一点展开，用 $0 \le x \le L$ 标记位置。式\eqref{eq:momentum-position}中的一阶导数仍给出动量在区间内部的作用。两端怎样相接，还需写入定义域。
判断一种边界条件是否给出自伴算符，要先求出它的伴随。

\begin{definition}{算符的伴随}{qm-adjoint}
    对稠密定义的$\op{A}$，若存在$\ket{\eta} \in \HH$使
    \begin{equation}
        \mel{\phi}{\op{A}}{\psi} = \braket{\eta}{\psi}
        \qquad\text{对所有 }\ket{\psi} \in \Dom(\op{A})
    \end{equation}
    成立，则$\ket{\phi}$属于$\Dom(\op{A}\adjoint)$，并定义
    $\op{A}\adjoint\ket{\phi} = \ket{\eta}$。
\end{definition}

等价地，要让 $\ket\phi$ 属于伴随的定义域，就要能找到常数 $C_\phi$，使每个允许的 $\ket\psi$ 都满足
$\abs{\mel{\phi}{\op A}{\psi}} \le C_\phi\norm{\psi}$。这时称 $\ket\psi \mapsto \mel{\phi}{\op A}{\psi}$ 为有界线性泛函(bounded linear functional)。
\begin{definition}{对称算符与自伴算符}{qm-self-adjoint}
    对稠密定义的 $\op A$，取任意 $\ket\phi,\ket\psi \in \Dom(\op A)$。对称性与自伴性的条件分别为
    \begin{subequations}
        \begin{align}
            \op A\text{ 对称}\quad &\Longleftrightarrow\quad
                \mel{\phi}{\op A}{\psi} = \mel{\psi}{\op A}{\phi}\conj,\\
            \op A\text{ 自伴}\quad &\Longleftrightarrow\quad
                \op A\text{ 对称且 }\Dom(\op A) = \Dom(\op A\adjoint).
        \end{align}
    \end{subequations}
    在有限维空间中，定义在整个空间上的 Hermite 矩阵满足 $\op{A} = \op{A}\adjoint$，因而自伴。
\end{definition}
\begin{remark}
    有些物理教材将对称(symmetric)也叫 Hermitian，本书在无穷维情形用“对称”避免混淆。
    对称性约束已有的本征值为实，却不足以保证完整的自伴谱理论和唯一的幺正时间演化。
    若对称算符的闭包自伴，就称原算符本质自伴(essentially self-adjoint)。此时这个闭包是它唯一的自伴扩张。
\end{remark}

现在比较区间 $[0,L]$ 上一阶导数的两种边界条件。对绝对连续、一阶导数平方可积的 $\phi,\psi$，分部积分给出
\begin{equation}\label{eq:qm-momentum-boundary-form}
    \int_0^L \phi\conj(-\i\hbar\psi')\d x
    - \int_0^L (-\i\hbar\phi')\conj\psi\d x
    = -\i\hbar[\phi\conj\psi]_0^L .
\end{equation}
先限定 $\op P$ 的定义域，要求态的位置波函数满足上述正则条件且在两端为零。取 $\ket\phi,\ket\psi$ 都在此域中，边界项消失，说明 $\op P$ 对称。
检验伴随时只要求 $\ket\psi$ 在此域中。取位置振幅恒为 $\phi(x) = 1/\sqrt L$ 的态 $\ket\phi$，它的振幅导数为零，边界项也因 $\psi(0) = \psi(L) = 0$ 消失。
于是对每个允许的 $\ket\psi$ 都有 $\mel{\phi}{\op P}{\psi} = 0$。由\dfnref{dfn:qm-adjoint}，$\op P\adjoint\ket\phi = 0$。由于 $\phi$ 的端点非零，$\ket\phi$ 超出了 $\op P$ 的定义域，这个 $\op P$ 因而不自伴\cite{domains}。
要让 $\op P\ket\psi$ 也属于 $\op P$ 的定义域，还须使 $\psi'$ 在两端为零。硬墙盒子的自伴动能算符要求二阶弱导数平方可积，以及波函数在两端为零，斜率可以非零。

\paragraph{环上的动量谱}
在环上要求 $\psi(L) = \e^{\i\theta}\psi(0)$，所有态共用同一个 $\theta$。先取支撑在区间内部的光滑 $\psi$，伴随的定义就要求 $\phi$ 有平方可积的弱导数，故在有限区间上可取绝对连续的代表。
再让 $\psi(0)$ 取任意值，式\eqref{eq:qm-momentum-boundary-form}右边对所有这样的 $\psi$ 都为零，恰好要求 $\phi(L) = \e^{\i\theta}\phi(0)$。
此时伴随与原算符有同一边界条件，都按 $-\i\hbar$ 乘以一阶导数作用，环上的动量因此自伴。

本征方程给出 $\braket{x}{p_n} = C\e^{\i p_nx/\hbar}$，代入端点条件得到 $\e^{\i p_n L/\hbar} = \e^{\i\theta}$。允许的动量为
\begin{equation}
    p_n = \frac\hbar L(2\pi n + \theta),\qquad n \in \ZZ.
\end{equation}

概率密度 $\abs{C}^2$ 为常数，在长为 $L$ 的环上归一化给出 $L\abs{C}^2 = 1$，所以 $\abs{C} = L^{-1/2}$。

半直线上的径向动量也受原点的边界条件限制。径向动量与自伴径向 Hamilton 算符的区别见\ssecref{ss:qm-radial-momentum}。

\subsectionstar{谱投影与算符函数}
\label{ss:qm-spectral-functions}
自旋的一个出口、环上的一个动量本征值、位置探测器的一组像素，都可以用“读数落在集合 $B$ 内”来描述。
记相应的投影为 $\op\Pi_A(B)$。全体读数的投影是 $\one$，互不相交的读数集合对应正交投影；
可数多个这样的集合取并集时，各投影之和作用于任一向量，按范数收敛到并集的投影。这样的投影族称为投影值测度(projection-valued measure, PVM)。

\begin{theorem}{自伴算符的谱定理}{qm-spectral-theorem}
    每个自伴算符 $\op{A}$ 唯一确定实轴上的投影值测度 $\op{\Pi}_A$，使 $\op{A}$ 在其定义域上等于谱积分
    \begin{equation}
        \op A = \int_{\RR}a\d\op\Pi_A(a).
    \end{equation}
\end{theorem}
对离散读数 $a$，单点集合的投影 $\op\Pi_A(\{a\})$ 就是 $\op\Pi_a$。若能用这些本征子空间展开整个态空间，谱积分便化为 $\textstyle \sum_a a\op\Pi_a$。
若有无限多个这样的本征值，写 $\op A\ket\psi$ 的展开时须有 $\ket\psi \in \Dom(\op A)$，各部分和按向量范数收敛。
对归一化态 $\ket\psi$，记 $\mu_\psi(B) = \mel{\psi}{\op\Pi_A(B)}{\psi}$，即 Born 规则给出的读数落在 $B$ 内的概率。
有限维时，对一个给定矩阵，所有单位向量经它作用后的长度都有共同的上界。在无限维空间，一个算符若没有这样的有限上界，就称为无界算符(unbounded operator)。
此时计算二阶矩可用 $\|\op A\ket\psi\|^2$。若写成矩阵元 $\mel{\psi}{\op A^2}{\psi}$，还要保证能对 $\op A\ket\psi$ 再作用一次 $\op A$。

如果把读数 $a$ 换算成 $f(a)$，可以用同一个谱投影构造相应的算符 $f(\op A)$。这个方法称为谱函数演算，适用于定义在 $\op A$ 的谱上的 Borel 可测标量函数 $f$。它也给出了新算符的定义域\cite{teschl-forms}，
\begin{subequations}\label{eq:qm-spectral-functional-calculus}
    \begin{gather}
        f(\op A) = \int_{\RR} f(a)\d\op\Pi_A(a),\\
        \Dom(f(\op A)) = \biggset{\ket\psi \in \HH}
            {\int_{\RR}\abs{f(a)}^2\d\mu_\psi(a) < \infty}.
    \end{gather}
\end{subequations}
离散本征子空间张成全空间时，第一式就是 $\textstyle\sum_a f(a)\op\Pi_a$，积分条件也相应化为对各本征值求和。实值 $f$ 给出自伴算符。若 $f$ 在谱上有界，$f(\op A)$ 就能作用于整个态空间。
如果 $\op\Pi_A(B) = 0$，改变 $f$ 在 $B$ 上的值不影响 $f(\op A)$，例如直线位置谱中的一个单点。

\paragraph{算符函数的例子}
对 $\op\sigma_x$ 取 $f(a) = a^2$，两个本征值 $\pm1$ 的平方都是 $1$，所以 $f(\op\sigma_x) = \one$。若直接把 $\op\sigma_x$ 在 $z$ 基中的矩阵元逐个平方，得到的却还是 $\op\sigma_x$ 的矩阵。

在位置表象中，$f(\op X)$ 将每个位置分量乘以 $f(x)$。$f$ 的定义域描述哪些\emph{数}可代入，$\Dom(f(\op X))$ 则描述哪些\emph{态矢}乘上 $f(x)$ 后仍平方可积。
以 $f(x) = 1/x$ 为例。直线上单点 $x = 0$ 不影响平方可积范数，可在该点任取一个有限值。前面的矩形波包属于 $\Dom(\op X)$，因为它有界支撑，但 $\psi(x)/x$ 在零点附近的模平方正比于 $1/x^2$，积分发散。因此这个波包超出了 $\Dom(f(\op X))$。

再取固定的实数 $a$，其量纲为长度，令 $f(p) = \e^{-\i ap/\hbar}$。$f(\op P)$ 在动量表象中将每个分量乘以 $f(p)$。这个乘数的模处处为 $1$，所以前面的矩形波包也属于 $f(\op P)$ 的定义域。
若把标量指数的 Taylor 级数逐项改写为 $\op P$ 的幂，一次项就需要 $\op P\ket\psi$，而矩形波包没有平方可积的一阶弱导数。第三章将用这个相位写出空间平移。

有些映射本来就以算符为输入。例如，有限维空间的迹把矩阵映为数，量子信道 $\mathcal E$ 把输入的密度算符映为新的密度算符（\dfnref{dfn:qm-cptp}）。这些映射各有自己的定义。

\paragraph{平均能量与能量算符的定义域}
在谱函数演算中取 $\op A = \op H$，就能解释\ssecref{ss:qm-regularity}的撤墙例子。先取能谱非负的自伴算符 $\op H$，如自由 Hamilton 算符 $\op H_0$。此时 $\mu_\psi$ 是能量读数的概率分布。若一阶矩有限，平均能量为
\begin{equation}
    \ave{\op H} = \int_{[0,\infty)} E\,\d\mu_\psi(E)
        = \norm{\op H^{1/2}\ket\psi}^2 .
    \label{eq:qm-energy-form-mean}
\end{equation}
最后一个等号由式\eqref{eq:qm-spectral-functional-calculus}取 $f(E) = \sqrt E$ 得到。$\op H^{1/2}$ 把能量为 $E$ 的振幅乘以 $\sqrt E$，求范数平方就给这一部分概率乘上 $E$。使一阶矩有限的态构成 $\Dom(\op H^{1/2})$，称为能量的\emph{二次型域}。若原能谱下界为 $b$，可先对 $\op H - b\one$ 用此式，再给均值加回 $b$。

要让 $\op H\ket\psi$ 本身成为向量，作用在振幅上的因子是 $E$，求范数平方时便得到 $E^2$。因此，算符的\emph{定义域}满足下列等价条件\cite{teschl-forms}。
\begin{equation}
    \ket\psi \in \Dom(\op H)
        \quad\Longleftrightarrow\quad
        \int_{[0,\infty)} E^2\,\d\mu_\psi(E) < \infty .
    \label{eq:qm-energy-operator-domain}
\end{equation}
在这个域内，$\mel{\psi}{\op H}{\psi}$ 就是 $\ket\psi$ 与 $\op H\ket\psi$ 的内积。
若只有一阶矩有限，平均能量由式\eqref{eq:qm-energy-form-mean}的二次型给出。也可以先只计入 $E \leq M$ 的读数对均值的贡献，再令 $M \to \infty$。

撤墙瞬间的波包 $\ket\phi$ 具有有限的平均动能 $\hbar^2\pi^2/(2mL^2)$，所以它属于 $\Dom(\op H_0^{1/2})$。将波函数延拓到整条实线后，$\phi''$ 含有端点的 $\delta$ 项，因此 $\ket\phi$ 超出了 $\Dom(\op H_0)$。由式\eqref{eq:qm-energy-operator-domain}，撤墙后能量的二阶矩和方差都发散。有墙时，测量这个基态的能量，每次都得到 $E_1$，方差为零。
这里的端点 $\delta$ 来自对初态求导，撤墙后的势能处处为零。\exmref{exm:qm-contact-well}考察含有 $g\delta(x)$ 势能项的系统，式\eqref{eq:qm-contact-slope-jump}给出它对斜率的要求。

\paragraph{算符函数的对易子}
设 $f$ 是光滑的位置函数。对光滑且有紧支撑的波包，$f(\op X)$ 乘以 $f(x)$，$\op P$ 求导。用乘积求导比较两种作用次序，
\begin{equation}\label{eq:qm-momentum-position-function}
    \mel{x}{[\op P,f(\op X)]}{\psi}
    = -\i\hbar\Bigp{\bigp{f(x)\psi(x)}' - f(x)\psi'(x)} = -\i\hbar f'(x)\psi(x).
\end{equation}
因此在这些波包上，$[\op P,f(\op X)] = -\i\hbar f'(\op X)$，其中 $f'(\op X)$ 由标量导数 $f'(x)$ 构造。取 $f = V$，第三章将用这个结果求平均力。

若 $f$ 是多项式，可以逐次展开乘积的对易子。涉及无界算符时，取一个共同的稠密定义域，要求 $\op A$ 和 $\op B$ 都把域中的向量映回该域。这样每个乘积都能在这个域上计算，并且
\begin{equation}
    [\op A,\op B^n]
    = \sum_{j=0}^{n-1}\op B^j[\op A,\op B]\op B^{n-1-j},\qquad n \geq 1 .
\end{equation}
若 $[[\op A,\op B],\op B] = 0$，每一项都等于 $\op B^{n-1}[\op A,\op B]$，于是对任意多项式 $f$ 有
$[\op A,f(\op B)] = f'(\op B)[\op A,\op B]$。
取 $\op A = \op P$、$\op B = \op X$，在 Schwartz 波包上便恢复式\eqref{eq:qm-momentum-position-function}对多项式 $f$ 的结果。
若两算符有界、$[[\op A,\op B],\op B] = 0$，且 $f$ 在原点的幂级数收敛半径大于 $\|\op B\|$，则对易子和导函数的级数都按算符范数收敛，公式也适用于这样的 $f$。

\reading{不确定关系与相容测量见 MIT 8.05 讲义 5\cite{mituncertainty}。区间上的边界条件怎样影响算符的自伴性和能谱，可参见 Bonneau 等\cite{domains}。下一章从给定的初态求出测量时刻的态，再计算各读数的概率。}
